\documentclass{amsart}
\usepackage{amsmath,amssymb,amsthm,hyperref}

\newtheorem{theorem}{Theorem}
\newtheorem{proposition}{Proposition}
\newtheorem{lemma}{Lemma}
\newtheorem{corollary}{Corollary}
\theoremstyle{definition}
\newtheorem{definition}{Definition}
\newtheorem{remark}{Remark}

\DeclareMathOperator{\E}{\mathbb{E}}
\newcommand{\R}{\mathbb{R}}
\newcommand{\N}{\mathbb{N}}
\newcommand{\one}{\mathbf{1}}

\begin{document}

\title{Resolving the zero-prior obstruction to confirming universal hypotheses}
\maketitle

\begin{abstract}
We give a complete and rigorous resolution of the zero-prior obstruction. We
prove that the obstruction is \emph{exactly} an artifact of restricting to
priors that are absolutely continuous with respect to Lebesgue measure on the
parameter space, and we characterize the class of priors for which the
universal hypothesis $H=\{\theta=1\}$ is confirmable. The admissible class is
the class of \emph{spike-and-slab} (mixed discrete--continuous) priors that put
a strictly positive atom on $\theta=1$. For every prior in this class we show
that (i) $P(H)>0$, (ii) the posterior probability of $H$ after $n$ consecutive
successes is strictly positive and increases to $1$ as $n\to\infty$, and (iii)
the framework is an ordinary Bayesian model (hence coherent under conditioning),
is invariant under reparametrizations of $\theta$, and reduces to the
Bayes--Laplace predictive rule in the non-universal regime. We then prove a
converse: an exchangeable model confirms $H$ in the sense of (ii) if and only if
its de~Finetti mixing measure has an atom at $\theta=1$. Finally we exhibit the
algorithmic-probability (Solomonoff) realization of the same construction and
show it satisfies (i)--(iii), explaining why a Solomonoff-style universal prior
escapes the obstruction automatically.
\end{abstract}

\section{Setup and the obstruction restated}\label{sec:setup}

\subsection{The model}
Fix the parameter space $\Theta=[0,1]$. For $\theta\in\Theta$ let $P_\theta$ be
the law on $\{0,1\}^\N$ under which $X_1,X_2,\dots$ are i.i.d.\ Bernoulli$(\theta)$;
thus for every $n$ and every $x_1,\dots,x_n\in\{0,1\}$,
\begin{equation}\label{eq:likelihood}
  P_\theta(X_1=x_1,\dots,X_n=x_n)=\theta^{\,s_n}(1-\theta)^{\,n-s_n},
  \qquad s_n:=\sum_{i=1}^n x_i .
\end{equation}
A prior is a Borel probability measure $\mu$ on $\Theta$. The Bayesian model it
generates is the measure
\begin{equation}\label{eq:joint}
  P(\,\cdot\,)=\int_{\Theta}P_\theta(\,\cdot\,)\,d\mu(\theta)
\end{equation}
on $\Theta\times\{0,1\}^\N$ (with $\theta$ the first coordinate). This is the
standard exchangeable Bayesian model; by de~Finetti's theorem every
exchangeable law on $\{0,1\}^\N$ arises this way for a unique $\mu$
(Theorem~\ref{thm:converse} below uses this).

The universal hypothesis is the event
\begin{equation}\label{eq:H}
  H=\{\theta=1\}\subseteq\Theta .
\end{equation}
We write $A_n=\{X_1=\dots=X_n=1\}$ for the event of $n$ consecutive successes,
and $A_\infty=\bigcap_{n\ge1}A_n=\{X_i=1\ \forall i\}$.

\subsection{The obstruction}
We restate the problem's premise precisely, to fix exactly what must be
circumvented.

\begin{proposition}[The zero-prior obstruction]\label{prop:obstruction}
If $\mu$ is absolutely continuous with respect to Lebesgue measure on $[0,1]$,
then $P(H)=0$ and $P(H\mid A_n)=0$ for every finite $n$.
\end{proposition}

\begin{proof}
If $\mu\ll\mathrm{Leb}$ then $\mu(\{1\})=0$ because $\{1\}$ is Lebesgue-null;
hence $P(H)=\mu(H)=\mu(\{1\})=0$. For the posterior, by~\eqref{eq:joint}
and~\eqref{eq:likelihood},
\[
  P(A_n)=\int_0^1\theta^{\,n}\,d\mu(\theta),\qquad
  P(H\cap A_n)=\int_{\{1\}}\theta^{\,n}\,d\mu(\theta)=1^n\cdot\mu(\{1\})=0 .
\]
Since $\mu\ll\mathrm{Leb}$ is not the point mass at $0$, $\mu$ charges some
subinterval of $(0,1]$, so $P(A_n)=\int_0^1\theta^n\,d\mu(\theta)>0$ and
conditioning is well defined; thus
$P(H\mid A_n)=P(H\cap A_n)/P(A_n)=0$.
\end{proof}

The proof isolates the sole mechanism of the obstruction: $\mu(\{1\})=0$. The
likelihood ratio in favour of $H$ over any $\theta<1$ after $n$ successes is
$1/\theta^n\to\infty$, so the data are maximally favourable to $H$; the only
reason the posterior of $H$ cannot rise is that the \emph{prior} assigns it
measure zero, and Bayesian updating can never lift an event of prior measure
zero. The cure is therefore forced: give $H$ positive prior mass. The content
of the theorems below is that doing so can be done while preserving every
desirable feature of the continuous model.

\section{The admissible class: spike-and-slab priors}\label{sec:class}

\begin{definition}[Spike-and-slab prior]\label{def:spikeslab}
A \emph{spike-and-slab prior} is a measure of the form
\begin{equation}\label{eq:prior}
  \mu = p\,\delta_{1} + (1-p)\,\nu,
\end{equation}
where $p\in(0,1)$, $\delta_1$ is the unit point mass at $\theta=1$, and $\nu$ is
a Borel probability measure on $[0,1]$ with $\nu(\{1\})=0$ (the ``slab'').
We call $p$ the \emph{spike weight} and $\nu$ the \emph{slab law}. We say the
slab is \emph{non-degenerate} if $\nu([1-\varepsilon,1))>0$ for every
$\varepsilon>0$, i.e.\ $1\in\operatorname{supp}\nu$; the canonical choice is
$\nu=\mathrm{Beta}(a,b)$ with $a,b>0$, and in particular the
Bayes--Laplace slab $\nu=\mathrm{Unif}[0,1]$.
\end{definition}

Throughout, $P$ denotes the model~\eqref{eq:joint} built from~\eqref{eq:prior}.
This is an ordinary prior (a probability measure on $[0,1]$), so the entire
machinery of Bayesian conditioning applies verbatim; in particular requirement
(iii)'s coherence is automatic. We record the elementary integrals once.

\begin{lemma}[Marginal likelihood of a success run]\label{lem:marginal}
For every $n\ge0$,
\begin{equation}\label{eq:Pan}
  P(A_n)=p+(1-p)\,m_n,\qquad m_n:=\int_{[0,1)}\theta^{\,n}\,d\nu(\theta),
\end{equation}
with the convention $A_0$ the whole space and $m_0=1$. Moreover $m_n$ is
non-increasing in $n$ and $\lim_{n\to\infty}m_n=0$.
\end{lemma}

\begin{proof}
By~\eqref{eq:joint}, \eqref{eq:likelihood} and~\eqref{eq:prior},
\[
  P(A_n)=\int_{[0,1]}\theta^n\,d\mu(\theta)
  =p\cdot 1^n+(1-p)\int_{[0,1]}\theta^n\,d\nu(\theta).
\]
Since $\nu(\{1\})=0$, the last integral equals $\int_{[0,1)}\theta^n\,d\nu=m_n$,
giving~\eqref{eq:Pan}. For $\theta\in[0,1)$ the sequence $\theta^n$ is
non-increasing in $n$, so $m_n$ is non-increasing by monotonicity of the
integral. Finally $\theta^n\to 0$ pointwise on $[0,1)$ and $0\le\theta^n\le1$,
so by dominated convergence (dominating function $\one$, which is
$\nu$-integrable) $m_n\to\int_{[0,1)}0\,d\nu=0$.
\end{proof}

\section{Confirmation: properties (i) and (ii)}\label{sec:confirm}

\begin{theorem}[Positive prior and convergent posterior]\label{thm:main}
Let $\mu=p\,\delta_1+(1-p)\,\nu$ be any spike-and-slab prior
(Definition~\ref{def:spikeslab}). Then:
\begin{enumerate}
\item[\textup{(i)}] $P(H)=p>0$.
\item[\textup{(ii)}] For every $n\ge0$ the posterior of $H$ given $n$ consecutive
successes is
\begin{equation}\label{eq:posterior}
  P(H\mid A_n)=\frac{p}{p+(1-p)\,m_n}\in(0,1],
\end{equation}
this quantity is non-decreasing in $n$, and
\[
  \lim_{n\to\infty}P(H\mid A_n)=1.
\]
\item[\textup{(iii)}] Conditioned on the full success sequence,
$P(H\mid A_\infty)=1$; i.e.\ the model gives the correct answer in the limit.
\end{enumerate}
\end{theorem}

\begin{proof}
(i) $P(H)=\mu(\{1\})=p\,\delta_1(\{1\})+(1-p)\,\nu(\{1\})=p+0=p>0$.

(ii) The event $A_n$ has $P(A_n)=p+(1-p)m_n\ge p>0$ by
Lemma~\ref{lem:marginal}, so the conditional probability is defined. The
contribution of $H$ to $A_n$ is the spike, whose likelihood is $1^n=1$:
\[
  P(H\cap A_n)=\int_{\{1\}}\theta^n\,d\mu(\theta)=p\cdot 1^n=p .
\]
Therefore
\[
  P(H\mid A_n)=\frac{P(H\cap A_n)}{P(A_n)}=\frac{p}{p+(1-p)\,m_n},
\]
which is~\eqref{eq:posterior}. Because $p>0$ and $(1-p)m_n\ge0$ the value lies
in $(0,1]$. The map $t\mapsto p/(p+(1-p)t)$ is strictly decreasing in
$t\ge0$, and $m_n$ is non-increasing in $n$ (Lemma~\ref{lem:marginal}); hence
$P(H\mid A_n)$ is non-decreasing in $n$. Finally, since $m_n\to0$,
\[
  \lim_{n\to\infty}P(H\mid A_n)=\frac{p}{p+(1-p)\cdot 0}=\frac{p}{p}=1 .
\]

(iii) The events $A_n$ decrease to $A_\infty$, and by~\eqref{eq:joint}
\[
  P(H\cap A_\infty)=\int_{\{1\}}P_\theta(A_\infty)\,d\mu(\theta)=p\cdot 1=p,
\]
because $P_1(A_\infty)=1$. Also
$P(A_\infty)=\int_{[0,1]}P_\theta(A_\infty)\,d\mu
=p\cdot1+(1-p)\int_{[0,1)}\big(\lim_n\theta^n\big)\,d\nu$, and
$\lim_n\theta^n=0$ for $\theta\in[0,1)$ while the point $\theta=1$ is
excluded by $\nu(\{1\})=0$; by dominated convergence the slab integral is $0$, so
$P(A_\infty)=p$. Hence $P(H\mid A_\infty)=p/p=1$. (Consistently, by downward
continuity of measure $P(H\mid A_\infty)=\lim_n P(H\mid A_n)=1$.)
\end{proof}

This settles (i) and (ii) for the entire admissible class, with no restriction
on $\nu$ beyond $\nu(\{1\})=0$: the slab may be uniform, any Beta law, or any
other continuous prior on $[0,1)$. In particular the spike weight $p$ may be
taken arbitrarily small, so the resolution does not require committing strong
prior belief in $H$.

\section{Coherence, reparametrization invariance, and predictive behaviour: (iii)}
\label{sec:coherence}

\subsection{Coherence under updating}
The model $P$ of~\eqref{eq:joint} is a genuine probability measure on
$\Theta\times\{0,1\}^\N$ and $H$ a genuine event with $P(H)=p>0$. All the
posteriors above are ordinary conditional probabilities of positive-probability
events, so Bayes coherence (the chain rule, conglomerability over the discrete
data, and sequential consistency: $P(H\mid A_{n+1})$ is the one-step update of
$P(H\mid A_n)$) holds automatically. Explicitly, the sequential update factor is
\begin{equation}\label{eq:sequential}
  \frac{P(H\mid A_{n+1})}{P(H\mid A_n)}
  =\frac{P(A_n)}{P(A_{n+1})}
  =\frac{p+(1-p)m_n}{p+(1-p)m_{n+1}}\ \ge 1,
\end{equation}
which is exactly Bayes' rule applied to the single new observation
$X_{n+1}=1$ (the inequality holds because $m_{n+1}\le m_n$), confirming that the
increasing posterior of Theorem~\ref{thm:main}(ii) is produced by lawful
conditioning and not by any extra-Bayesian device.

\subsection{Invariance under reparametrization}
A reparametrization is a relabelling of the parameter by a Borel isomorphism
$\phi:[0,1]\to[0,1]$; we write $\eta=\phi(\theta)$ for the new coordinate and
$\psi:=\phi^{-1}$ for its inverse, so $\theta=\psi(\eta)$. A reparametrization
must leave the actual statistical model unchanged: the trial that has success
probability $\theta$ in the old chart has success probability $\psi(\eta)$ in
the new chart, so the likelihood of a success run becomes
$\eta\mapsto\psi(\eta)^n$, and the prior is transported by push-forward,
$\tilde\mu:=\phi_\#\mu$. The hypothesis $H=\{\theta=1\}$ becomes
$\tilde H:=\phi(H)=\{\eta=\phi(1)\}$. We require $\phi(1)=1$ (equivalently
$\psi(1)=1$): the value ``success probability $=1$'', which \emph{defines} the
universal hypothesis, is identified across charts.

\begin{proposition}[Reparametrization invariance]\label{prop:invariance}
Let $\phi:[0,1]\to[0,1]$ be a Borel isomorphism with $\phi(1)=1$, let
$\psi=\phi^{-1}$, and let $\tilde P$ be the model in the new chart: prior
$\tilde\mu=\phi_\#\mu$ and likelihood $\tilde P_\eta(A_n)=\psi(\eta)^n$.
\begin{enumerate}
\item[\textup{(a)}] \textup{(Class is closed.)} If $\mu=p\delta_1+(1-p)\nu$ is a
spike-and-slab prior then $\tilde\mu=p\,\delta_1+(1-p)\,\tilde\nu$ with
$\tilde\nu=\phi_\#\nu$, and $\tilde\nu(\{1\})=\nu(\psi(\{1\}))=\nu(\{1\})=0$;
so $\tilde\mu$ is again spike-and-slab with the same spike weight $p$.
\item[\textup{(b)}] \textup{(Posterior is invariant.)} For every $n$,
\[
  \tilde P(\tilde H\mid A_n)=P(H\mid A_n).
\]
\end{enumerate}
\end{proposition}

\begin{proof}
(a) Push-forward of a sum of measures is the sum of push-forwards, and
$\phi_\#\delta_1=\delta_{\phi(1)}=\delta_1$. Since $\phi$ is a bijection with
$\phi(1)=1$ we have $\psi(\{1\})=\{1\}$, so
$\tilde\nu(\{1\})=(\phi_\#\nu)(\{1\})=\nu(\psi(\{1\}))=\nu(\{1\})=0$.

(b) We compute every quantity in the new chart and invoke the change-of-variables
(image-measure) formula $\int g\,d(\phi_\#\mu)=\int (g\circ\phi)\,d\mu$.
The new marginal of a success run is
\[
  \tilde P(A_n)=\int_{[0,1]}\psi(\eta)^n\,d\tilde\mu(\eta)
  =\int_{[0,1]}\psi(\eta)^n\,d(\phi_\#\mu)(\eta)
  =\int_{[0,1]}\psi(\phi(\theta))^n\,d\mu(\theta)
  =\int_{[0,1]}\theta^n\,d\mu(\theta)=P(A_n),
\]
using $\psi\circ\phi=\mathrm{id}$. Likewise the new joint mass of the universal
hypothesis is, since $\tilde H=\{\eta=1\}$ and $\psi(1)=1$,
\[
  \tilde P(\tilde H\cap A_n)
  =\int_{\{1\}}\psi(\eta)^n\,d\tilde\mu(\eta)
  =\psi(1)^n\,\tilde\mu(\{1\})
  =1^n\cdot p=p=P(H\cap A_n).
\]
Dividing, $\tilde P(\tilde H\mid A_n)=p/P(A_n)=P(H\mid A_n)$.
\end{proof}

\begin{remark}[What invariance does and does not mean]\label{rem:invariance}
Proposition~\ref{prop:invariance} shows the framework is genuinely
coordinate-free: both the admissible class and the entire posterior trajectory
$\{P(H\mid A_n)\}_n$ are invariant under any relabelling of $\theta$ that fixes
the success-endpoint. This is exactly the structural feature the obstruction
\emph{lacked} for continuous priors. Absolute continuity is \emph{not} a
coordinate-free property capable of charging $H$: no Borel isomorphism can turn a
Lebesgue-null point into a positive-Lebesgue-measure set, so a prior with
$\mu\ll\mathrm{Leb}$ can never be transported into one with $\tilde\mu(\{1\})>0$.
By contrast an \emph{atom} at the success-endpoint is intrinsic: a point mass
pushes forward to a point mass of the same weight, and the image of the
success-endpoint under any admissible relabelling is again the success-endpoint.
Hence ``$P(H)>0$'' is invariant under reparametrization, whereas
``$\mu\ll\mathrm{Leb}$'' is the chart-dependent constraint responsible for the
obstruction.
\end{remark}

\subsection{Predictive behaviour in the non-universal regime}
We verify that away from the pure-success path the model behaves like an
ordinary, well-calibrated continuous Bayesian learner, and in particular
recovers the Bayes--Laplace rule of succession.

\begin{proposition}[Reduction to standard predictive behaviour]\label{prop:predictive}
Use the spike-and-slab prior~\eqref{eq:prior} with a slab $\nu$ that charges
$(0,1)$.
\begin{enumerate}
\item[\textup{(a)}] \emph{(Falsification.)} If at any stage some $X_i=0$ is
observed, then on that event the posterior of $H$ is $0$ and stays $0$ forever;
formally $P(H\mid X_1=x_1,\dots,X_n=x_n)=0$ whenever $s_n<n$.
\item[\textup{(b)}] \emph{(Continuous learning off the spike.)} Conditional on
any data string containing a $0$, the posterior over $\theta$ is the ordinary
Bayesian posterior of the slab model $\nu$; the spike contributes nothing.
\item[\textup{(c)}] \emph{(Rule of succession; recovery of Bayes--Laplace.)}
With the uniform slab $\nu=\mathrm{Unif}[0,1]$ one has $m_n=\frac1{n+1}$ and
\[
  P(X_{n+1}=1\mid A_n)=\frac{p+(1-p)\frac{1}{n+2}}{\,p+(1-p)\frac{1}{n+1}\,}.
\]
As $p\downarrow0$ this equals $\tfrac{n+1}{n+2}$, Laplace's rule of succession,
for every $n$; for fixed $p\in(0,1)$ it lies strictly above $\tfrac{n+1}{n+2}$
and increases to $1$, the (mild) optimism induced by the spike.
\end{enumerate}
\end{proposition}

\begin{proof}
(a) If $s_n<n$ then $x_j=0$ for some $j\le n$, and the spike's likelihood is
$P_1(X_1=x_1,\dots,X_n=x_n)=1^{s_n}(1-1)^{n-s_n}=0$ because $n-s_n\ge1$. Thus
$P(H\cap\{X_1=x_1,\dots,X_n=x_n\})=p\cdot0=0$, while the slab gives the data
positive probability because $\nu$ charges $(0,1)$, where
$\theta^{s_n}(1-\theta)^{n-s_n}>0$; so the posterior of $H$ is
$0/(\text{positive})=0$. As $H$ has posterior $0$, all further conditioning keeps
it at $0$ (an event of conditional probability $0$ cannot be revived).

(b) On data with $s_n<n$ the spike has zero likelihood by (a), so in the
posterior decomposition
$d\mu_{\text{post}}(\theta)\propto\theta^{s_n}(1-\theta)^{n-s_n}\,d\mu(\theta)$
the atom at $1$ receives weight $0$ and the surviving mass is
$\propto\theta^{s_n}(1-\theta)^{n-s_n}\,d\nu(\theta)$, which is precisely the
posterior obtained from the slab prior $\nu$ alone. Hence all predictions
coincide with those of the continuous model $\nu$.

(c) For $\nu=\mathrm{Unif}[0,1]$, $m_n=\int_0^1\theta^n\,d\theta=\frac1{n+1}$.
By the chain rule and Lemma~\ref{lem:marginal},
\[
  P(X_{n+1}=1\mid A_n)=\frac{P(A_{n+1})}{P(A_n)}
  =\frac{p+(1-p)\frac{1}{n+2}}{p+(1-p)\frac{1}{n+1}} .
\]
Setting $p=0$ gives $\frac{1/(n+2)}{1/(n+1)}=\frac{n+1}{n+2}$, Laplace's rule.
For $p\in(0,1)$: writing the predictive value as $f(p)$, one has $f(0)=\frac{n+1}{n+2}$
and $f(1)=1$; since $\frac1{n+1}>\frac1{n+2}>0$, cross-multiplication shows
$p+(1-p)\frac1{n+2}$ and $p+(1-p)\frac1{n+1}$ are both increasing in $p$ but the
ratio $f(p)$ is strictly increasing in $p$ on $[0,1]$ (its derivative has the sign
of $\frac1{n+1}-\frac1{n+2}>0$), so $f(p)>\frac{n+1}{n+2}$ for $p\in(0,1)$; and
as $n\to\infty$ both numerator and denominator tend to $p$, so $f(p)\to1$.
\end{proof}

Parts (a)--(b) deliver the philosophically essential feature: the spike is a
\emph{falsifiable} hypothesis in the Popperian sense---a single counterexample
($X_i=0$) annihilates its posterior---yet it accrues confirmation under a
growing run of confirming instances. This is exactly the asymmetry one wants for
``all ravens are black'': one non-black raven refutes it; many black ravens
progressively confirm it.

\section{Characterization: a converse}\label{sec:converse}

We now show that the spike is not merely \emph{sufficient} but \emph{necessary}:
within the entire class of exchangeable Bernoulli models, confirmability of $H$
in the sense of Theorem~\ref{thm:main}(ii) is equivalent to the presence of an
atom at $\theta=1$. This is what makes the admissible class a genuine
\emph{characterization}.

\begin{theorem}[Confirmability $\iff$ atom at the success-endpoint]\label{thm:converse}
Let $\mu$ be any prior on $[0,1]$, with associated exchangeable model $P$
of~\eqref{eq:joint}, and suppose $\mu\ne\delta_0$ so that $P(A_n)>0$ for all $n$.
Set $p:=\mu(\{1\})$. Then for all $n$
\begin{equation}\label{eq:general}
  P(H\mid A_n)=\frac{p}{\int_{[0,1]}\theta^n\,d\mu(\theta)} ,
\end{equation}
and consequently:
\begin{enumerate}
\item[\textup{(a)}] $P(H\mid A_n)>0$ for some (equivalently every) $n$ $\iff$
$p>0$;
\item[\textup{(b)}] $\displaystyle\lim_{n\to\infty}P(H\mid A_n)=1 \iff p>0$.
\end{enumerate}
In particular, an exchangeable Bernoulli model confirms the universal hypothesis
$H$ (in the sense of property (ii): posterior strictly positive and converging to
$1$ along a success run) \emph{if and only if} its de~Finetti mixing measure
$\mu$ places strictly positive mass on $\theta=1$. Every such $\mu$ with $p<1$
is a spike-and-slab prior (Definition~\ref{def:spikeslab}) with slab
$\nu=(\mu-p\delta_1)/(1-p)$, and $\mu=\delta_1$ when $p=1$.
\end{theorem}

\begin{proof}
By~\eqref{eq:likelihood}, $P(A_n)=\int_{[0,1]}\theta^n\,d\mu(\theta)$, and
$P(H\cap A_n)=\int_{\{1\}}\theta^n\,d\mu=p$, which gives~\eqref{eq:general}; the
denominator is positive since $\mu\ne\delta_0$ implies $\mu$ charges $(0,1]$,
where $\theta^n>0$.

(a) If $p=0$ then the numerator of~\eqref{eq:general} is $0$, so
$P(H\mid A_n)=0$ for every $n$. If $p>0$ the numerator is positive and the
denominator finite and positive, so $P(H\mid A_n)>0$ for every $n$. (The
``some/every'' equivalence holds because the numerator does not depend on $n$.)

(b) If $p=0$ then $P(H\mid A_n)\equiv0\not\to1$. Suppose $p>0$. If $p=1$ then
$\mu=\delta_1$, $P(A_n)=1$ and $P(H\mid A_n)=1$ for all $n$. If $0<p<1$, write
$\mu=p\,\delta_1+(1-p)\nu$ with $\nu=(\mu-p\delta_1)/(1-p)$, a probability
measure with $\nu(\{1\})=(\mu(\{1\})-p)/(1-p)=0$. Then
$\int_{[0,1]}\theta^n\,d\mu=p+(1-p)m_n$ with $m_n=\int_{[0,1)}\theta^n\,d\nu\to0$
by dominated convergence exactly as in Lemma~\ref{lem:marginal}. Hence
$P(H\mid A_n)=p/(p+(1-p)m_n)\to p/p=1$. The displayed iff statements combine (a)
and (b).

The identification of $\mu$ with the unique mixing measure of an exchangeable law
is the classical de~Finetti representation theorem for $\{0,1\}$-valued
exchangeable sequences; uniqueness of $\mu$ makes ``$\mu(\{1\})>0$'' a
well-defined, model-intrinsic criterion, independent of any parametrization.
\end{proof}

Theorem~\ref{thm:converse} is the precise sense in which the spike-and-slab
class is \emph{the} answer: no exchangeable Bayesian model outside this class can
confirm $H$, and every model inside it does, with the convergence rate governed
entirely by the slab through $m_n$ (faster mass concentration of $\nu$ near $1$
gives slower confirmation, and vice versa). The continuous Bayes--Laplace prior
fails precisely and only because it is the degenerate case $p=0$.

\section{The Solomonoff / algorithmic-probability realization}\label{sec:solomonoff}

We now show that the same resolution arises canonically, without an \emph{ad hoc}
choice of spike weight, when one replaces the parametric prior by a universal
(Solomonoff) prior on hypotheses; this answers the algorithmic-probability part
of the problem and explains \emph{why} algorithmic induction is immune to the
obstruction.

\subsection{Construction}
Work directly on the space of \emph{hypotheses} rather than on $[0,1]$. Let
$\mathcal M=\{Q_1,Q_2,\dots\}$ be a countable class of computable measures on
$\{0,1\}^\N$, chosen to contain
\begin{equation}\label{eq:Qstar}
  Q_\star:=\delta_{(1,1,1,\dots)}\quad\text{(the deterministic ``all successes'' law)},
\end{equation}
i.e.\ $Q_\star(A_n)=1$ for all $n$ and $Q_\star(\{X_i=0\})=0$; this $Q_\star$ is
the algorithmic incarnation of the universal hypothesis $H$. Fix a prefix-free
universal description and let $K(\cdot)$ denote prefix Kolmogorov complexity;
define the universal mixture (a Solomonoff-style prior over $\mathcal M$)
\begin{equation}\label{eq:Mprior}
  M(\cdot)=\sum_{k\ge1} w_k\,Q_k(\cdot),\qquad
  w_k:=\frac{2^{-K(Q_k)}}{Z},\quad
  Z:=\sum_{j\ge1}2^{-K(Q_j)}\in(0,1] .
\end{equation}
Each $w_k>0$ because $K(Q_k)<\infty$ (every $Q_k$ is computable, hence has a
finite description), and $\sum_k w_k=1$. Let $w_\star:=w_{k_\star}>0$ be the
weight of $Q_\star$ in~\eqref{eq:Mprior}.

The Solomonoff posterior on the hypothesis $Q_\star$ is the ordinary Bayesian
posterior of the discrete prior $(w_k)_k$:
\begin{equation}\label{eq:solpost}
  M(Q_\star\mid A_n)=\frac{w_\star\,Q_\star(A_n)}{\sum_k w_k\,Q_k(A_n)}
  =\frac{w_\star}{\sum_k w_k\,Q_k(A_n)} .
\end{equation}

\begin{theorem}[Solomonoff confirmation of the universal hypothesis]\label{thm:sol}
With $M$ as in~\eqref{eq:Mprior} and $Q_\star$ as in~\eqref{eq:Qstar}:
\begin{enumerate}
\item[\textup{(i)}] $M$ assigns the universal hypothesis strictly positive prior
weight: the prior probability of $Q_\star$ is $w_\star>0$.
\item[\textup{(ii)}] $M(Q_\star\mid A_n)\in(0,1]$ is non-decreasing in $n$. It
strictly exceeds $w_\star$ once some $Q_k$ with positive weight has $Q_k(A_n)<1$,
and
$\lim_{n\to\infty}M(Q_\star\mid A_n)
=\dfrac{w_\star}{\sum_k w_k\,Q_k(A_\infty)}$;
this limit equals $1$ if and only if $Q_\star$ is the only law in $\mathcal M$
with $Q_k(A_\infty)>0$ (e.g.\ when every other $Q_k$ produces a $0$ eventually,
$Q_k$-almost surely).
\item[\textup{(iii)}] The scheme is coherent (it is plain Bayesian conditioning
of the discrete prior $(w_k)$), it is invariant under the choice of universal
machine up to a multiplicative $O(1)$ factor in the weights (hence does not
affect the limit in (ii)), and in the non-universal regime it reduces to the
ordinary Solomonoff predictor, which predicts consistently.
\end{enumerate}
\end{theorem}

\begin{proof}
(i) The prior probability of $Q_\star$ under~\eqref{eq:Mprior} is by definition
$w_\star=2^{-K(Q_\star)}/Z>0$, since $Q_\star=\delta_{(1,1,\dots)}$ is computable
($K(Q_\star)<\infty$).

(ii) From~\eqref{eq:solpost}, $Q_\star(A_n)=1$, and for each $k$,
$Q_k(A_n)\in[0,1]$ is non-increasing in $n$ (because $A_{n+1}\subseteq A_n$).
Hence the denominator $D_n:=\sum_k w_k Q_k(A_n)$ is non-increasing in $n$ and
bounded below by $w_\star Q_\star(A_n)=w_\star>0$, so
$M(Q_\star\mid A_n)=w_\star/D_n$ is non-decreasing and lies in $(0,1]$; it
exceeds $w_\star=M(Q_\star)$ as soon as $D_n<1$, which occurs once some
positively-weighted $Q_k$ has $Q_k(A_n)<1$. By monotone convergence
(equivalently dominated convergence with summable dominating sequence $(w_k)$),
\[
  D_n=\sum_k w_k Q_k(A_n)\ \longrightarrow\ \sum_k w_k\,Q_k(A_\infty),
\]
since $Q_k(A_n)\downarrow Q_k(A_\infty)$ for each $k$. Therefore
$M(Q_\star\mid A_n)\to w_\star/\sum_k w_k Q_k(A_\infty)$. As
$Q_\star(A_\infty)=1$, this limit is $1$ iff
$\sum_{k:Q_k\ne Q_\star} w_k Q_k(A_\infty)=0$, i.e.\ iff every other $Q_k$ has
$Q_k(A_\infty)=0$; this is the algorithmic analogue of the condition
$\nu(\{1\})=0$ in Theorem~\ref{thm:converse}.

(iii) Coherence holds because~\eqref{eq:solpost} is Bayes' rule for the discrete
prior $(w_k)$. Machine-independence is the Invariance Theorem of Kolmogorov
complexity: for two universal prefix machines $U,U'$ there is a constant
$c=c(U,U')$ with $|K_U(Q)-K_{U'}(Q)|\le c$ for all $Q$; hence the unnormalized
weights satisfy $2^{-c}\le 2^{-K_U(Q)}/2^{-K_{U'}(Q)}\le 2^{c}$, a bounded
multiplicative perturbation that leaves $w_\star>0$ and (being a fixed positive
factor up to normalization) does not change the limit
$M(Q_\star\mid A_n)\to1$ in (ii). Finally, on data containing a $0$ the term
$Q_\star(A_n)=0$ kills $Q_\star$'s posterior, and $M$ reverts to the standard
Solomonoff mixture over $\mathcal M\setminus\{Q_\star\}$, whose predictive
consistency (Solomonoff's theorem: $M(X_{n+1}=1\mid X_{1:n})$ converges, with
probability one, to the true conditional probability for any computable data
source) gives reasonable behaviour in the non-universal regime.
\end{proof}

\subsection{Why algorithmic probability resolves the obstruction
``for free''}
The parametric model required us to insert the spike by hand because the
continuum $[0,1]$ contains the universal hypothesis as a measure-zero point and
the only translation-/scale-natural continuous measures (Lebesgue and its
absolutely continuous cousins) cannot charge it. In the algorithmic setting the
hypothesis space is the \emph{countable} class of computable laws, on which the
universal prior~\eqref{eq:Mprior} is necessarily \emph{discrete}; \emph{every}
computable hypothesis, including the deterministic $Q_\star$, automatically
receives positive weight $\propto2^{-K(Q)}$. Thus the algorithmic prior is a
spike-and-slab prior in disguise---a discrete mixture with an atom on the
universal law---and the obstruction simply never arises. The complexity weighting
$2^{-K(Q_\star)}$ additionally supplies a non-arbitrary value for the spike
weight: simple universal regularities (such as ``always $1$'') are short to
describe, hence are \emph{a priori} preferred, which is precisely Occam's razor
operating in favour of confirmable universal generalizations.

\section{Summary: adequacy of the resolution}\label{sec:summary}

We collect the resolution and its justification.

\begin{enumerate}
\item \textbf{The obstruction is exactly absolute continuity}
(Proposition~\ref{prop:obstruction}): it depends only on $\mu(\{1\})=0$, a
chart-dependent constraint, not on any structural feature of confirmation.

\item \textbf{Admissible class.} The spike-and-slab priors
$\mu=p\delta_1+(1-p)\nu$, $p\in(0,1)$, $\nu(\{1\})=0$
(Definition~\ref{def:spikeslab}), give $P(H)=p>0$ and posterior
$P(H\mid A_n)=p/(p+(1-p)m_n)\uparrow1$
(Theorem~\ref{thm:main}); this is requirements (i)--(ii).

\item \textbf{Coherence and invariance} (requirement (iii)):
the model is an ordinary Bayesian model, so updating is coherent
(\eqref{eq:sequential}); the admissible class is closed under reparametrization
and the entire posterior trajectory is invariant
(Proposition~\ref{prop:invariance}, Remark~\ref{rem:invariance}); and the model
reduces to the Bayes--Laplace rule of succession in the non-universal regime,
with one observed counterexample falsifying $H$ outright
(Proposition~\ref{prop:predictive}).

\item \textbf{Characterization} (Theorem~\ref{thm:converse}): among
\emph{all} exchangeable Bernoulli models, confirmability of $H$ holds if and only
if the de~Finetti mixing measure has an atom at $\theta=1$. The class of item~2
is therefore not one device among many but the complete solution set.

\item \textbf{Algorithmic realization} (Theorem~\ref{thm:sol}): a
Solomonoff-style universal prior over computable hypotheses is automatically a
discrete mixture charging the deterministic ``all successes'' law $Q_\star$ with
weight $2^{-K(Q_\star)}>0$; it confirms $Q_\star$ with posterior tending to $1$,
is machine-independent up to $O(1)$, and reduces to the consistent Solomonoff
predictor off the success path. Here the spike weight is fixed canonically by
Occam's razor rather than chosen by hand.
\end{enumerate}

\noindent\textbf{Philosophical adequacy.} The construction reconciles three
desiderata usually thought to be in tension. (a) \emph{Confirmability}: $H$ is
genuinely learnable, with monotone increasing posterior tending to certainty
under accumulating positive instances. (b) \emph{Falsifiability}: a single
counterexample drives the posterior of $H$ to $0$ irreversibly
(Proposition~\ref{prop:predictive}(a)), so $H$ remains an empirical, refutable
claim and the framework is not credulous. (c) \emph{Open-mindedness}: the slab
$\nu$ keeps positive belief on every $\theta<1$, so before decisive evidence the
model entertains the possibility that $H$ is false and predicts in a calibrated,
Laplace-like manner. The zero-prior obstruction is thus revealed as an avoidable
artifact of an over-restrictive (merely-continuous) prior class; once the
genuinely universal hypothesis is admitted as a first-class atom---whether by an
explicit spike or, canonically, through algorithmic probability---standard
Bayesian conditioning confirms it exactly as scientific practice demands.
\end{document}
